{
    \setlength{\parindent}{0em}
    \par {\zihao{4}\bfseries 一、题目：\Title}
    \\
    \par {\zihao{4}\bfseries 二、指导教师对文献综述、开题报告、外文翻译的具体要求：}
}

本课题围绕多模态视频大模型高效推理关键技术开展研究，重点探索将视觉 token 压缩方法迁移应用于视频大模型结构中的可行性与有效性。

开题报告应在充分梳理多模态大模型与视频大模型发展现状的基础上，明确课题研究的必要性与现实意义。特别应突出视频模型中视觉 token 数量快速增长所带来的推理效率瓶颈，并结合相关文献分析提出具有针对性的技术路线。开题报告字数要求达到3500字以上。
同时，围绕多模态大模型高效推理完成10-15篇英文文献的阅读，并完成3000字以上的文献综述，要求阐明本领域的国内外发展现状、主要研究内容、进展情况、存在问题和参考依据。

此外，希望通过阅读文献和撰写开题报告阶段，提高解决实际问题的能力， 养成严谨、细致的工作作风，为后续科研工作打下良好基础。

\mbox{} \vfill

\signature{指导教师（签名）}
% comment the line above and uncomment the line below if you want to set a signature with a specific date.
% \signaturewithdate{指导教师（签名）}{1897}{5}{21}
